% \title{Apunte de Matemáticas Discretas CC3101:\\ 
% Caminos más cortos  Clase 20}
% % \title{Ejercicios Lógica e Inducción}
% \author{Andrés Abeliuk}
% \date{}
% \maketitle

\section{Caminos más cortos sobre grafos}

\subsection{Motivación}
En muchas aplicaciones prácticas, como el diseño de redes de transporte o de comunicación, necesitamos encontrar el camino más corto entre dos puntos. Este problema puede modelarse usando grafos con pesos, donde los vértices representan ubicaciones y las aristas representan las conexiones con sus respectivos costos o distancias.

\section{Grafos con pesos}
\begin{definitionblock}{Grafo con pesos}
Un \textbf{grafo con pesos} es una terna $(V, E, w)$ donde:
\begin{itemize}
  \item $V$ es el conjunto de vértices,
  \item $E \subseteq V \times V$ es el conjunto de aristas,
  \item $w: E \rightarrow \mathbb{R}_{\geq 0}$ es una función que asigna un peso a cada arista.
\end{itemize}
\end{definitionblock}

Sea $P = (v_0, v_1, \dots, v_k)$ un camino en un grafo con pesos. Su \textbf{longitud} es la suma de los pesos de sus aristas, es decir:
\[
\ell(P) = \sum_{i=0}^{k-1} w(v_i, v_{i+1})
\]

Un \textbf{camino más corto} entre dos vértices $u$ y $v$ es aquel camino $P$ que parte en $u$, termina en $v$, y tiene longitud mínima entre todos los posibles caminos entre ellos. Este camino no necesariamente es único.




\begin{exampleblock}{Ejemplo}
La siguiente red representa distancias entre ciudades:

\begin{center}
\begin{tikzpicture}[every node/.style={scale=0.8}]
  \tikzstyle{LabelStyle}=[sloped, fill=white]
  \Vertex[x=0,y=4]{a}
  \Vertex[x=2,y=4]{b}
  \Vertex[x=4,y=4]{c}
  \Vertex[x=1,y=2.5]{d}
  \Vertex[x=3,y=2.5]{e}
  \Vertex[x=2,y=1]{f}
  \Vertex[x=4,y=1]{g}

  \draw (a) to node[midway,fill=white,sloped] {$3$} (b);
  \draw (b) to node[midway,fill=white,sloped] {$2$} (c);
  \draw (a) to node[midway,fill=white,sloped] {$6$} (d);
  \draw (b) to node[midway,fill=white,sloped] {$2$} (d);
  \draw (c) to node[midway,fill=white,sloped] {$6$} (e);
  \draw (d) to node[midway,fill=white,sloped] {$3$} (f);
  \draw (e) to node[midway,fill=white,sloped] {$2$} (f);
  \draw (f) to node[midway,fill=white,sloped] {$1$} (g);
  \draw (c) to node[midway,fill=white,sloped] {$3$} (g);
\end{tikzpicture}
\end{center}

\vspace{0.5em}
\textbf{Pregunta:} ¿Cuál es la ruta más corta desde $a$ hasta $g$?
\end{exampleblock}




\subsection{Algoritmo de Dijkstra}
El algoritmo de Dijkstra es uno de los más importantes en teoría de grafos y fue propuesto por Edsger W. Dijkstra en 1956. Su objetivo es encontrar el camino más corto desde un vértice de origen hasta todos los demás vértices de un grafo con pesos no negativos.

El procedimiento sigue una estrategia codiciosa (greedy): comienza con el nodo de origen y explora iterativamente el vértice adyacente más cercano, acumulando las distancias mínimas conocidas. A cada paso, actualiza las distancias si encuentra un camino más corto que los previamente registrados.

Este algoritmo ha sido fundamental en el desarrollo de sistemas de navegación, planificación de rutas y optimización de redes, y es una pieza clave en muchos cursos de algoritmos y ciencias de la computación teórica.

\subsection*{Ejemplo}
Aplicaremos el algoritmo de Dijkstra desde el vértice $a$ para encontrar el camino más corto hacia todos los demás vértices.

\begin{center}
\begin{tikzpicture}[every node/.style={scale=0.8}]
  \tikzstyle{LabelStyle}=[sloped, fill=white]
  \Vertex[x=0,y=4]{a}
  \Vertex[x=2,y=4]{b}
  \Vertex[x=4,y=4]{c}
  \Vertex[x=1,y=2.5]{d}
  \Vertex[x=3,y=2.5]{e}
  \Vertex[x=2,y=1]{f}
  \Vertex[x=4,y=1]{g}

  \draw[very thick, red] (a) to node[midway,fill=white,sloped] {$3$} (b);
  \draw[very thick, red] (b) to node[midway,fill=white,sloped] {$2$} (c);
  \draw (a) to node[midway,fill=white,sloped] {$6$} (d);
  \draw (b) to node[midway,fill=white,sloped] {$2$} (d);
  \draw (c) to node[midway,fill=white,sloped] {$6$} (e);
  \draw (d) to node[midway,fill=white,sloped] {$3$} (f);
  \draw (e) to node[midway,fill=white,sloped] {$2$} (f);
  \draw (f) to node[midway,fill=white,sloped] {$1$} (g);
  \draw[very thick, red] (c) to node[midway,fill=white,sloped] {$3$} (g);
\end{tikzpicture}
\end{center}

\medskip
\noindent\textbf{Resolución del algoritmo de Dijkstra:}


\begin{center}
\renewcommand{\arraystretch}{1.3}
\begin{tabular}{|c|ccccccc|c|}
\hline
\textbf{Iteración} & \textbf{a} & \textbf{b} & \textbf{c} & \textbf{d} & \textbf{e} & \textbf{f} & \textbf{g} & \textbf{Procesado} \\
\hline
1 & $0_a$ & $3_a$ & $\infty$ & $6_a$ & $\infty$ & $\infty$ & $\infty$ & $a$ \\
2 & $0_a$ & $3_a$ & $5_b$ & $5_b$ & $\infty$ & $\infty$ & $\infty$ & $b$ \\
3 & $0_a$ & $3_a$ & $5_b$ & $5_b$ & $11_c$ & $\infty$ & $8_c$ & $c$ \\
4 & $0_a$ & $3_a$ & $5_b$ & $5_b$ & $11_c$ & $8_d$ & $8_c$ & $d$ \\
5 & $0_a$ & $3_a$ & $5_b$ & $5_b$ & $11_c$ & $8_d$ & $8_c$ & $f$ \\
6 & $0_a$ & $3_a$ & $5_b$ & $5_b$ & $10_f$ & $8_d$ & $8_c$ & $e$ \\
\hline
\end{tabular}
\end{center}

\medskip
\noindent\textbf{Camino más corto desde $a$ hasta $g$}: $a \rightarrow b \rightarrow c \rightarrow g$, con costo total $8$.



\begin{itemize}
  \item La celda $x_y$ indica la distancia más corta encontrada hasta el nodo $x$, con $y$ indicando el nodo anterior en el camino más corto.
  \item Se termina cuando todos los nodos están procesados.
  \item El camino más corto a un nodo se puede reconstruir siguiendo los subíndices hacia atrás.
\end{itemize}

\subsubsection{Pseudocódigo comentado}
El siguiente pseudocódigo describe el algoritmo de Dijkstra para encontrar caminos más cortos desde un nodo origen hacia todos los demás nodos de un grafo con pesos no negativos. 

\begin{verbatim}
function Dijkstra(Graph, source):
  // Inicialización: asigna distancia infinita a todos los nodos y sin predecesor conocido
  for each vertex v in Graph:
    dist[v] := inf              // Distancia estimada desde el origen a v
    prev[v] := undefined        // Nodo anterior en el camino más corto encontrado

  dist[source] := 0             // La distancia del nodo origen a sí mismo es 0
  Q := all nodes in Graph       // Conjunto de nodos no visitados

  while Q is not empty:
    u := node in Q with min dist[u]   // Selecciona el nodo más cercano aún no procesado
    remove u from Q
    
    for each neighbor v of u:
      alt := dist[u] + length(u, v)   // Calcula la distancia alternativa a través de u
      if alt < dist[v]:
        dist[v] := alt                // Se encontró un camino más corto hacia v
        prev[v] := u                  // u se convierte en el predecesor de v

  return dist[], prev[]               // dist contiene las distancias mínimas; prev los caminos
\end{verbatim}


\subsection{Observaciones importantes}
El algoritmo de Dijkstra es versátil y puede aplicarse tanto en grafos no dirigidos como en \textbf{grafos dirigidos}, lo cual amplía significativamente su utilidad práctica en modelar distintos tipos de redes. 

\textbf{ Restricción de pesos no negativos:} Es importante destacar que este algoritmo requiere que todos los pesos de las aristas sean no negativos, ya que su lógica de actualización de distancias mínimas depende de esta propiedad. En presencia de pesos negativos, el algoritmo puede fallar al no detectar caminos más cortos correctamente. 

\textbf{Tratamiento de múltiples aristas:} Si el grafo tiene múltiples aristas entre un mismo par de nodos, se recomienda conservar únicamente aquella con menor peso, ya que siempre será preferible en la búsqueda de caminos mínimos. 

% \textbf{3. Aplicabilidad general:} Gracias a su diseño, Dijkstra puede ser implementado eficientemente en distintas estructuras de datos, y su uso se ha extendido ampliamente en sistemas de navegación, planificación de rutas y redes computacionales.



% \section{Ejercicio propuesto}
% \begin{thmblock}{Aplicar Dijkstra desde $A$ hacia $D$}
% \begin{minipage}{0.45\textwidth}
% % \includegraphics[width=0.9\textwidth]{DijsktraExampleGraph.png}
% \end{minipage}\hfill
% \begin{minipage}{0.45\textwidth}
% % \includegraphics[width=0.9\textwidth]{DijsktraExampleTable.png}
% \end{minipage}
% \end{thmblock}

\subsection*{Más recursos}
    
Para ver el algoritmo en funcionamiento, sobre un grafo ingresado por el usuario, ver \textbf{demo interactiva:} \
\url{https://algorithms.discrete.ma.tum.de/graph-algorithms/spp-dijkstra/index_en.html}